% =====================================================================
%  Conferenze — Introduzione alla parte
%  Presenta i due eventi (GAMM di Roma e Zagarolo) e il criterio con cui
%  le trascrizioni sono state trasformate in capitoli.
%  Capitolo non numerato. Nessun pacchetto aggiuntivo.
% =====================================================================

\chapter*{Introduzione: due conferenze}
\phantomsection
\addcontentsline{toc}{chapter}{Introduzione: due conferenze}

Questa parte raccoglie gli interventi seguiti di persona a \textbf{due
conferenze} sul videogioco, entrambe a Roma e dintorni:
\begin{enumerate}
  \item il \textbf{Checkpoint Festival}, alla sua \textbf{terza edizione}, ospitato
        dal \textbf{GAMM} (Game Museum) di Roma: cinque incontri su videogiochi e
        società, sviluppo indipendente, preservazione e musica (capitoli C1.x);
  \item \textbf{On The Road -- III Conferenza internazionale sul Game Design}, a
        \textbf{Zagarolo} (Roma), su due giornate: sei interventi il primo giorno
        (capitoli C2.x) e cinque il secondo (capitoli C3.x), in italiano e in
        inglese, dedicati soprattutto all'industria vista da dentro: publishing,
        produzione, carriera, eventi, tecnologia.
\end{enumerate}

\section*{Come leggere questi capitoli}

Ogni capitolo corrisponde a un \textbf{talk}, un'\textbf{intervista} o un
\textbf{panel}, ed è ricostruito a partire da una trascrizione automatica
dell'audio. Valgono quindi alcune convenzioni:
\begin{itemize}
  \item il contenuto segue l'ordine dell'intervento e cerca di non omettere nulla di
        significativo, comprese le domande del pubblico;
  \item gli interventi in inglese sono resi in italiano; le citazioni tra virgolette
        sono traduzioni fedeli nel senso, non alla lettera;
  \item i \textbf{nomi propri} (persone, studi, giochi) sono spesso storpiati dalla
        trascrizione: quando la forma corretta è ricostruibile la si usa,
        segnalando in nota la forma trascritta; quando non lo è, lo si dichiara;
  \item le parti \textbf{incomprensibili} e le ripetizioni meccaniche della
        trascrizione sono segnalate in nota e omesse;
  \item quando il genere o il nome completo di una persona non sono noti, si usa
        una formulazione neutra (``l'ospite'', ``il relatore'').
\end{itemize}
Ogni capitolo si chiude, come nel resto del libro, con i \textbf{punti chiave}.
